% Part: first-order-logic
% Chapter: models-theories
% Section: expressing-props-of-structures

\documentclass[../../../include/open-logic-section]{subfiles}

\begin{document}

\olfileid{fol}{mat}{exs}

\olsection{\printtoken{P}{structure}-এর ধর্ম প্রকাশ}

\begin{explain}
অপেক্ষক ও সম্পর্কের উপর আরোপিত শর্ত, অথবা আরও সাধারণভাবে কোনো গঠনের
অপেক্ষক ও সম্পর্কগুলি ওই শর্তগুলি পরিতৃপ্ত করে—এ কথা প্রকাশ করা প্রায়ই
উপকারী ও গুরুত্বপূর্ণ। যেমন, কোনো ভাষার যেসব !!{structure}s আমাদের
কাঙ্ক্ষিত অর্থে !!{predicate}s-গুলিকে “ধারণ” করে, সেগুলিকে অন্য
গঠনগুলি থেকে আলাদা করার উপায় চাই। অবশ্য কোন !!{structure}s-কে
আমরা “অভিপ্রেত” বলব তা নির্দিষ্ট করার সম্পূর্ণ স্বাধীনতা আমাদের আছে।
যেমন, আমরা বলতে পারি যে !!{predicate}~$\le$-এর ব্যাখ্যা একটি ক্রম হতেই
হবে; অথবা আমরা শুধু~$\Lang L$-এর সেই ব্যাখ্যাগুলিতে আগ্রহী যাদের
সংজ্ঞাক্ষেত্র সেট দিয়ে গঠিত এবং যেখানে $\Obj \in$-কে
“!!a{element} হওয়া” সম্পর্ক দিয়ে ব্যাখ্যা করা হয়। কিন্তু
!!a{structure}-এর ভাষার !!{sentence}s দিয়েই কি আমরা এ কথা বলতে পারি?
অন্যভাবে বললে,~$\Struct M$-এর ভাষার কোনো !!a{sentence} (অথবা সম্ভবত
!!{sentence}s-এর কোনো সেট) দিয়ে~$\Struct M$-এর উপর আরোপিত কোন
শর্তগুলি প্রকাশ করতে পারি? কিছু শর্ত আমরা প্রকাশ করতে পারব না। যেমন,
ভাষা~$\Lang L_A$-এর এমন কোনো বাক্য নেই যা কোনো
!!{structure}~$\Struct M$-এ কেবল তখনই সত্য, যখন $\Domain M = \Nat$।
“সংজ্ঞাক্ষেত্রে শুধু স্বাভাবিক সংখ্যাই আছে”—এ কথা প্রকাশ করা যায় না।
কিন্তু !!{structure}s-এর কিছু “গঠনগত ধর্ম” সম্ভবত প্রকাশ করা যায়।
!!{sentence}s দিয়ে !!{structure}s-এর কোন ধর্মগুলি প্রকাশ করতে পারি?
অথবা অন্যভাবে বললে, !!a{sentence} (বা !!{sentence}s-এর সেট)-কে সত্য
করা !!{structure}s-এর সংগ্রহ হিসেবে কোন সংগ্রহগুলিকে বর্ণনা করতে পারি?
\end{explain}

\begin{defn}[কোনো সেটের মডেল]
ধরি, $\Gamma$ ভাষা~$\Lang L$-এর !!{sentence}s-এর একটি সেট।
প্রত্যেক $!A \in \Gamma$-র জন্য $\Sat{M}{!A}$ হলে বলি,
!!a{structure}~$\Struct M$ \emph{$\Gamma$-র একটি মডেল}।
\end{defn}

\begin{ex}
বাক্য~$\lforall[x][x \le x]$ গঠন~$\Struct M$-এ সত্য যদি এবং কেবল যদি
$\Assign{\le}{M}$ একটি প্রতিবিম্ব সম্পর্ক। বাক্য
$\lforall[x][\lforall[y][((x \le y \land y \le x) \lif x = y)]]$
গঠন~$\Struct M$-এ সত্য যদি এবং কেবল যদি $\Assign{\le}{M}$ একটি
প্রতিসম-বিরোধী সম্পর্ক। বাক্য
$\lforall[x][\lforall[y][\lforall[z][((x \le y \land y \le z)
      \lif x \le z)]]]$ গঠন~$\Struct M$-এ সত্য যদি এবং কেবল যদি
$\Assign{\le}{M}$ একটি পরিযায়ী সম্পর্ক। তাই
\begin{align*}
\{\quad &\lforall[x][x \le x], \\
   & \lforall[x][\lforall[y][((x \le y \land y \le
    x) \lif x = y)]], \\
   &\lforall[x][\lforall[y][\lforall[z][((x \le y
      \land y \le z) \lif x \le z)]]] \quad \}
\end{align*}
-এর মডেলগুলি ঠিক সেইসব গঠন যেখানে~$\Assign{\le}{M}$ প্রতিবিম্ব,
প্রতিসম-বিরোধী ও পরিযায়ী, অর্থাৎ একটি আংশিক ক্রম। ফলে এগুলিকে
\emph{আংশিক ক্রমের প্রথম-ক্রমের তত্ত্বের} স্বতঃসিদ্ধ হিসেবে নিতে পারি।
\end{ex}

\end{document}
